\section{E4: explicit density calibration}
\label{sec:density}

\lead{The prediction.} The first report measured a direct trade-off: every increase in
layer-1 density costs layer-1 accuracy, and the specificity parameter $s$ is a blunt way to
choose a point on it. If density is what layer~2 needs, then setting it directly --- rather
than hoping $s$ lands somewhere usable --- should be strictly better, and there should be a
target rate at which the stack works.

\lead{The controller does its job.} Table~\ref{tab:calib} shows it reaching every target it is
given, from the greedy layer's own \code{0.14\%} up by two orders of magnitude, by removing
literals from clauses that fire too rarely and adding them to clauses that fire too often. It
costs seconds, it leaves no channel dead, and it is the same layer~1 in every row --- the same
cached automata from E1, adjusted.

\begin{table}[h]
\centering\small
\caption{The target sweep. \code{target} is the requested per-patch firing rate, in percent;
\code{raw} and \code{pooled} are the median rates achieved before and after the
$\poolK\times\poolK$ OR pool; \code{size} is the median clause size the controller left behind
(the uncalibrated greedy layer has \densityGreedySize{}); \code{in band} the fraction of
channels landing inside the tolerance; \code{rounds} the number of single-literal adjustments
needed. Accuracy is the full two-layer stack, three seeds.}
\label{tab:calib}
\pgfplotstabletypeset[
  col sep=comma, string type,
  every head row/.style={before row=\ttoprule, after row=\ttmidrule},
  every last row/.style={after row=\ttbotrule},
  columns={target,fire_raw,fire_pooled,size,in_band,rounds,acc,sd},
  columns/target/.style={column name=\thead{target (\%)}, column type={r}},
  columns/fire_raw/.style={column name=\thead{raw (\%)}, column type={r}},
  columns/fire_pooled/.style={column name=\thead{pooled (\%)}, column type={r}},
  columns/size/.style={column name=\thead{size}, column type={r}},
  columns/in_band/.style={column name=\thead{in band}, column type={r}},
  columns/rounds/.style={column name=\thead{rounds}, column type={r}},
  columns/acc/.style={column name=\thead{test acc (\%)}, column type={r}},
  columns/sd/.style={column name=\thead{$\pm$}, column type={r}},
]{data/calib_sweep.csv}
\end{table}

\begin{figure}[h]
\centering
\begin{tikzpicture}
\begin{axis}[
  width=0.92\textwidth, height=6.2cm,
  xlabel={target layer-1 firing rate (\%, log scale)}, ylabel={test accuracy (\%)},
  xmode=log, log basis x=10,
  grid=major, grid style={ttline!40},
  legend style={font=\sffamily\scriptsize, draw=ttline, at={(0.5,-0.24)}, anchor=north,
                legend columns=3, column sep=1.2em},
  ymin=12, ymax=41,
  label style={font=\sffamily\small}, tick label style={font=\sffamily\footnotesize},
  every axis plot/.append style={very thick},
]
\addplot[ttorange, mark=*, mark size=2pt, error bars/.cd, y dir=both, y explicit]
  table[col sep=comma, x=target, y=acc, y error=sd] {data/calib_sweep.csv};
\addplot[ttslate, dashed, domain=0.3:30, samples=2] {\accMctmRandom};
\addplot[ttmoss, dotted, domain=0.3:30, samples=2] {\accMctm};
\legend{calibrated greedy layer 1, random layer 1 (uncalibrated), greedy layer 1 (uncalibrated)}
\end{axis}
\end{tikzpicture}
\caption{Stack accuracy against the layer-1 firing rate the controller was asked for. Density
buys a large, monotone improvement and levels off well short of the random-clause control.}
\label{fig:calib}
\end{figure}

\subsection{Density is a real effect, and it is not the whole effect}

\begin{ttkey}
Raising the target rate raises stack accuracy monotonically, from
\accMctmCalibRZeroZeroFive{}\% at a $0.5\%$ target to \accMctmCalibRTwoZero{}\% at $20\%$ ---
a gain of nearly ten points over the uncalibrated greedy stack's \accMctm{}\%, from a
post-hoc adjustment that costs seconds and trains nothing. Criterion~1 is met comfortably.

Criterion~2 is not, and the way it fails is the result. At the $20\%$ target the calibrated
layer fires on \calibFireMctmCalibRTwoZero{}\% of pooled positions --- \emph{more} than the
\fireRandomPre{}\% of the random-clause control --- and still reaches only
\accMctmCalibRTwoZero{}\% against \accMctmRandom{}\%. At matched density, and then past it,
the gap does not close.
\end{ttkey}

So density was a real cause and not the only one. The residual is clause \emph{size}, and the
sweep makes it visible without needing a separate experiment: the controller reaches
\calibFireMctmCalibRTwoZero{}\% by cutting the greedy clause from \densityGreedySize\ literals to
\calibSizeMctmCalibRTwoZero, while a random clause reaches a comparable rate with three. Both fire about equally often; one is a
\calibSizeMctmCalibRTwoZero-term description of a specific patch that happens to be common,
the other a three-term description of a general property. They have the same density and very different generality, and layer~2
can only build on the second.

Calibrating the random layer the other way confirms it from the other side. Taking it from
\fireRandomPre{}\% to the same $20\%$ target --- which for three-literal clauses means
\emph{removing} literals, not adding them --- moves it from \accMctmRandom{}\% to
\accMctmRandomCalib{}\%.

Two corollaries follow, both of which matter more than the headline number.

\lead{Density is necessary and weakly sufficient.} Every arm at a low firing rate is bad; not
every arm at a high firing rate is good. A controller that fixes only the rate leaves the
representation's other degeneracy --- long, brittle conjunctions --- untouched. A second
target, on clause size, is the obvious extension, and \S\ref{sec:sizetarget} runs it.

\lead{The controller is a general tool, not a patch.} Because it writes automata states one
step from the include boundary, it composes with continued training. \S\ref{sec:credit} uses
it that way, as a stabiliser inside a training loop rather than a one-off prune, and it is
what keeps layer~1 alive there at all.

\subsection{Testing the residual: a second target, on clause size}
\label{sec:sizetarget}

If size is the residual, the controller should be asked for it. Walking a clause towards a
size is the same loop with a different stop condition, so the only real question is
\emph{which} literal to remove: under a size cap, dropping the most load-bearing literal
leaves the most permissive ones behind and sends the firing rate towards~1, dropping the least
load-bearing leaves the most selective and sends it towards~0. Neither is the clause being
asked for. Removing the literal that leaves the rate closest to target is, and the
post-removal rate is available in closed form --- removing literal $k$ lets through exactly
the patches it alone was blocking --- so the greedy step is exact.

\begin{table}[h]
\centering\small
\caption{The size sweep, all at the $20\%$ firing-rate target that won Table~\ref{tab:calib}.
\code{target} is the requested maximum clause size and \code{size} what the controller
achieved; \code{raw} and \code{pooled} are the resulting firing rates in percent, \code{in
band} the fraction of channels still inside the rate tolerance. The uncalibrated greedy layer
has \densityGreedySize\ literals; the random control has three.}
\label{tab:size}
\pgfplotstabletypeset[
  col sep=comma, string type,
  every head row/.style={before row=\ttoprule, after row=\ttmidrule},
  every last row/.style={after row=\ttbotrule},
  columns={target,size,fire_raw,fire_pooled,dead,in_band,acc,sd},
  columns/target/.style={column name=\thead{target}, column type={r}},
  columns/size/.style={column name=\thead{size}, column type={r}},
  columns/fire_raw/.style={column name=\thead{raw}, column type={r}},
  columns/fire_pooled/.style={column name=\thead{pooled}, column type={r}},
  columns/dead/.style={column name=\thead{dead}, column type={r}},
  columns/in_band/.style={column name=\thead{in band}, column type={r}},
  columns/acc/.style={column name=\thead{test acc (\%)}, column type={r}},
  columns/sd/.style={column name=\thead{$\pm$}, column type={r}},
]{data/size_sweep.csv}
\end{table}

\begin{ttkey}
Cutting the clauses harder is worth more than the firing rate alone. Accuracy rises as the cap
falls --- \accMctmSizeKTwoFour{}\% at 24 literals, \accMctmSizeKOneTwo{}\% at 12,
\accMctmSizeKSix{}\% at 6 --- against \accMctmCalibRTwoZero{}\% for the same rate target with
no cap, and then turns over: at 3 literals the clause can no longer hold the requested rate
(only \inbandMctmSizeKThree\ of channels stay inside the band, against 1.00 everywhere above
it) and accuracy falls back to \accMctmSizeKThree{}\%. The optimum is interior, at
\sizeBestTarget\ literals and \sizeBestAcc\,$\pm$\,\sizeBestSd{}\%, and it is the first arm with a
\emph{trained} first layer, in either report, to beat the \accSingleMatched{}\% a single layer
reaches at the same total clause budget. It does not reach the \accSingleRf{}\% of a single
layer given the stack's $9\times9$ receptive field.
\end{ttkey}

\lead{The rate target does not determine the size, and the size is what matters.} The
size-capped rows are both shorter \emph{and} denser than the rate-only row, so the two are
worth disentangling. The control row does it. \code{mctm-calib-r20-tight} asks for the same
$20\%$ rate with the tolerance band narrowed from $2\times$ to $1.2\times$, so the walk cannot
stop at the band edge and has to keep cutting until it lands on the target. It ends at
\calibSizeMctmCalibRTwoZeroTight\ literals and \accMctmCalibRTwoZeroTight{}\%, so roughly half
of the size-capped arms' advantage is indeed just ``cut further''. The other half is not: a
\calibSizeMctmCalibRTwoZeroTight-literal clause and a \calibSizeMctmSizeKSix-literal one both
satisfy the rate band, and the short one is worth \sizeBestAcc{}\% against the long one's
\accMctmCalibRTwoZeroTight{}\%.

A firing-rate target therefore leaves a range of clause sizes feasible, and the short end of
that range is where the accuracy is. Asking for the rate alone stops as soon as the rate is
met, which at \calibSizeMctmCalibRTwoZeroTight\ literals it already is; asking for a size as
well is what pushes the walk past that point. The practical conclusion is blunt: delete all
but a handful of literals from every clause of a greedily trained Tsetlin feature layer,
choosing what to keep by firing rate.

\lead{What it does not do is beat an untrained layer.} \accMctmRandom{}\% still stands, from
three random literals per clause and no training at all. The best repaired layer here gets
within a few points of it after \epochsRun\ epochs of supervised layer-1 training followed by
surgery that removes nine tenths of what that training produced. It is hard to read that as
anything other than a statement about the value of the training.

\subsection{The same controller on an untrained layer 1}
\label{sec:randcalib}

Every arm so far starts from a layer~1 that was trained to classify, because that is what the
proposal specifies. \S\ref{sec:sizetarget} ends by observing that the most effective thing to
do to such a layer is to delete most of it. The obvious next question is what happens if it is
never trained in the first place --- the controller needs no labels, and the random-clause
control of the first report is already the strongest first layer in either report.

\begin{table}[h]
\centering\small
\caption{The firing-rate sweep run over a \emph{random} layer~1 instead of a trained one.
Columns as in Table~\ref{tab:calib}: \code{target} is the requested per-patch rate in percent,
\code{size} the median clause size the controller settles on, \code{raw} and \code{pooled} the
achieved rates. The uncalibrated random layer has three literals per clause and fires on
\fireRandomPre{}\% of pooled positions.}
\label{tab:randcalib}
\pgfplotstabletypeset[
  col sep=comma, string type,
  every head row/.style={before row=\ttoprule, after row=\ttmidrule},
  every last row/.style={after row=\ttbotrule},
  columns={target,size,fire_raw,fire_pooled,in_band,acc,sd},
  columns/target/.style={column name=\thead{target (\%)}, column type={r}},
  columns/size/.style={column name=\thead{size}, column type={r}},
  columns/fire_raw/.style={column name=\thead{raw (\%)}, column type={r}},
  columns/fire_pooled/.style={column name=\thead{pooled (\%)}, column type={r}},
  columns/in_band/.style={column name=\thead{in band}, column type={r}},
  columns/acc/.style={column name=\thead{test acc (\%)}, column type={r}},
  columns/sd/.style={column name=\thead{$\pm$}, column type={r}},
]{data/rand_sweep.csv}
\end{table}

\begin{figure}[h]
\centering
\begin{tikzpicture}
\begin{axis}[
  width=0.92\textwidth, height=6.4cm,
  xlabel={target layer-1 firing rate (\%, log scale)}, ylabel={test accuracy (\%)},
  xmode=log, log basis x=10, ymin=12, ymax=45,
  grid=major, grid style={ttline!40},
  legend style={font=\sffamily\scriptsize, draw=ttline, at={(0.5,-0.26)}, anchor=north,
                legend columns=3, column sep=1.0em},
  label style={font=\sffamily\small}, tick label style={font=\sffamily\footnotesize},
  every axis plot/.append style={very thick},
]
\addplot[ttorange, mark=*, mark size=2pt, error bars/.cd, y dir=both, y explicit]
  table[col sep=comma, x=target, y=acc, y error=sd] {data/rand_sweep.csv};
\addplot[ttslate, mark=square*, mark size=2pt, error bars/.cd, y dir=both, y explicit]
  table[col sep=comma, x=target, y=acc, y error=sd] {data/calib_sweep.csv};
\addplot[ttmoss, dashed, domain=0.3:70, samples=2] {\accSingleRf};
\legend{random layer 1, greedy layer 1, best single layer}
\end{axis}
\end{tikzpicture}
\caption{The same controller, the same targets, over a randomly initialised first layer and
over a greedily trained one. Neither curve is flat, which is the density effect; the gap
between them is what \epochsRun\ epochs of supervised layer-1 training cost. The dashed line
is the strongest single-layer baseline, \accSingleRf{}\%; the \emph{un}calibrated random layer
sits within a fifth of a point of it, at \accMctmRandom{}\%, and is left off for legibility.}
\label{fig:randcalib}
\end{figure}

\begin{ttkey}
This is the configuration that works. The curve has an interior optimum --- accuracy climbs
from \accMctmRcRZeroFive{}\% at a $5\%$ target to \randBestAcc{}\% at \randBestTarget{}\% and
falls back to \accMctmRcRSixZero{}\% at $60\%$, where the clauses are so general they carry
almost nothing --- and at its peak the stack reaches
\randBestAcc\,$\pm$\,\randBestSd{}\%, against \accSingleMatched{}\% for a single layer at the
same clause budget and \accSingleRf{}\% for a single layer given the stack's receptive field.
It clears all three criteria of \S\ref{sec:criteria}, and it is the first time in either
report that the two-layer stack has been worth its clause budget on CIFAR-10.

Layer~1 is \emph{never trained}. It is random clauses, adjusted by a label-free controller that
does nothing but count how often each one fires.
\end{ttkey}

The layer-2 clauses tell the same story as \S\ref{sec:l2learned} does on \code{near}: over
this first layer a layer-2 clause carries \sizeMctmRandomCalib\ literals with
\posMctmRandomCalib\ positive and \negMctmRandomCalib{}\% negations, against \sizeMctm\
literals with \posMctm\ positive and \negMctm{}\% negations over the greedy layer. It is still
not a short rule, but it is an order of magnitude closer to one. The size-target arms of
\S\ref{sec:sizetarget} go further still --- \negMctmSizeKThree{}\% negations at a three-literal
first layer --- so on CIFAR-10 too, the arms whose second layer stops asserting absences are
the arms that work.

Read alongside \S\ref{sec:sizetarget}, the two sweeps say the same thing from opposite
directions. Cutting a trained layer down to \sizeBestTarget\ literals recovers most of the gap;
starting from clauses that are already that short and never training them at all closes it.
Whatever \epochsRun\ epochs of supervised layer-1 training contribute on this task, it is
negative, and the controller is what makes an untrained layer usable rather than merely
better than a trained one.

\subsection{And when combined with \ideaA}

There is almost nothing for the controller to do. The autoencoder layer already fires at
\fireAutoPre{}\% with \sizeAutoPre-literal clauses; asked for the same $20\%$ target, the
controller moves it to \calibFireMctmAutoCalib{}\% and leaves the clause size at
\calibSizeMctmAutoCalib, and the arm gives \accMctmAutoCalib{}\% against \accMctmAuto{}\%
uncalibrated. That is the cleanest statement of what \ideaA\ buys on CIFAR-10: an unsupervised
objective puts the layer where the controller would have put it anyway, without being told to,
and the two do not compound because they are doing the same job.
